\chapter{Conclusion and Recommendations}\label{chap:chapter7}

%%%%%%%%%%%%%%%%%%%%%%%%%%%%%%%%%%%%%%%%%%%%%%%%%%
\section{Conclusion}\label{sec:conclusion}
%%%%%%%%%%%%%%%%%%%%%%%%%%%%%%%%%%%%%%%%%%%%%%%%%%

Socius is a dual-agent, AI-powered social communication chatbot designed to
support Filipino emerging adults in developing pragmatic competence through
culturally grounded, realistic conversational roleplay. Grounded in Grice's
Cooperative Principle, Havighurst's Developmental Tasks, Arnett's theory of
Emerging Adulthood, and Filipino cultural values---including \textit{Hiya},
\textit{Pakikisama}, \textit{Utang na Loob}, \textit{Kapwa}, and
\textit{Paggalang}---the system provides users with immersive scenario-based
practice and structured, data-driven feedback on their communicative performance.

This chapter evaluates the extent to which Socius fulfilled each of its four
research objectives, based on findings from functional testing, expert validation,
and end-user testing.

\subsection*{Objective 1: Dual-Agent Architecture}

\textit{To design and implement a dual-agent architecture (AI Partner and AI Mentor)
utilizing the Google Gemini API to facilitate a natural roleplay experience and
deep pragmatic analysis.}

\medskip

This objective was substantially achieved. End-to-end integration testing
confirmed that the full APAM Framework---from scenario loading and AI Partner
roleplay, through conversation logging, to AI Mentor analysis and feedback
generation---was implemented and functioned as a cohesive pipeline. The Holistic
Mentor agent additionally demonstrated the capacity to aggregate longitudinal user
data and recommend contextually appropriate scenarios, extending the system
beyond static practice into adaptive, personalized coaching.

However, end-user testing revealed a notable weakness in the naturalness of the
AI Partner's roleplay. Users reported that AI Partner responses were often too long
and occasionally repetitive, which disrupted conversational immersion. Although
the system achieved technical integration, the qualitative dimension of natural,
fluid roleplay remains an area requiring targeted prompt engineering and output
constraint refinement in future iterations.

\subsection*{Objective 2: Theoretical Framework for Scenario Generation}

\textit{To formulate a theoretical framework for scenario generation that synthesizes
Grice's Cooperative Principle, Developmental Psychology (Havighurst and Arnett),
and Filipino Cultural Values.}

\medskip

This objective was fully achieved. The synthesized theoretical framework---
referred to as the APAM Framework---was successfully operationalized into a
curated set of conversational scenarios that reflect the social developmental tasks
of Filipino emerging adults aged 18 to 24. Expert validation confirmed that the
framework is structurally sound, contextually nuanced, and capable of generating
socially realistic scenarios within the local cultural context. The scenarios
themselves were further validated by the expert, who affirmed their alignment with
the theoretical constructs underpinning the system.

User feedback corroborated this finding qualitatively. Participants consistently
provided positive remarks regarding the cultural relevance, diversity, and realism
of the scenarios, with several noting that conversations felt authentic enough to
elicit genuine emotional responses. This affirms that the framework achieved its
intended function of grounding social practice in scenarios that are both
developmentally appropriate and culturally resonant for the target population.

\subsection*{Objective 3: Multi-Level Quantitative Scoring Rubric}

\textit{To develop a multi-level, quantitative scoring rubric for pragmatic
competence---enforced via JSON Schema Controlled Generation---that provides
detailed, actionable metrics across Intention Coherence, Social \& Affective
Nuance, and Communicative Effectiveness.}

\medskip

This objective was achieved at both the design and technical levels. Expert
validation confirmed that the three-domain, eight-subcategory scoring rubric
constitutes a valid and theoretically grounded approach to quantifying pragmatic
ability. The use of JSON Schema Controlled Generation---enforcing structured,
float-scored outputs from the AI Mentor (Pragmatic Analysis) agent---was confirmed
to be technically sound through functional testing, ensuring that scores are
machine-parsable, reliably stored in the database, and consistently applied across
sessions.

In practice, individual users received personalized pragmatic scores per scenario
session, and the AI Mentor's feedback was structured and guided by the schema
outputs. Longitudinal averaging of scores across sessions further enabled
rudimentary tracking of user progress over time. Nonetheless, because the scoring
is generated by a probabilistic LLM without concurrent human validation, the
consistency and accuracy of individual score assignments remain subject to variance
inherent in the model.

\subsection*{Objective 4: Evaluation of Chatbot Effectiveness}

\textit{To evaluate the chatbot's effectiveness in enhancing users' pragmatic
competence and social communication skills through user feedback and human user
testing.}

\medskip

This objective was achieved as an evaluation process; however, the results
present a nuanced picture. End-user testing was successfully conducted with
participants engaging across multiple scenario sessions over a one-week period.
Qualitatively, user perceptions were broadly favorable: participants reported that
Socius provided valuable, personalized feedback and that the experience of
navigating realistic social scenarios contributed to their awareness of pragmatic
dynamics in Filipino interpersonal communication.

Quantitatively, however, no statistically significant changes were observed in
users' pragmatic ability scores across the testing period. This outcome is
consistent with established findings in behavioral and social science research,
where observable competence development typically requires sustained exposure
over weeks or months rather than days. As such, the absence of significant
quantitative change does not negate the system's promise; rather, it underscores
the insufficiency of the one-week testing window as a means of capturing genuine
pragmatic growth, and highlights the need for longer longitudinal evaluation in
future studies.

%%%%%%%%%%%%%%%%%%%%%%%%%%%%%%%%%%%%%%%%%%%%%%%%%%
\section{Limitations}\label{sec:limitations}
%%%%%%%%%%%%%%%%%%%%%%%%%%%%%%%%%%%%%%%%%%%%%%%%%%

\subsection{Technical Limitations}\label{subsec:tech-limitations}

The technical performance of Socius was subject to constraints arising from both
the infrastructure and the underlying language model, and these two sources of
limitation should be understood as distinct but interacting bottlenecks.

At the infrastructure level, the system was hosted on Railway's free tier, which
imposed restrictions on compute availability, uptime, and request handling
capacity. Expert evaluators and user testers encountered intermittent latency and
occasional service interruptions during sessions---particularly during the
computationally intensive AI Mentor (Pragmatic Analysis) calls, which were
configured with a thinking budget of 8,000 tokens. These disruptions affected
session continuity and user immersion, and may have influenced the depth of
engagement that participants were able to sustain during testing. Technical
reliability was accordingly identified as requiring urgent remediation, and
represents a deployment-level constraint that is separable from the system's
architectural quality.

At the model level, the system's reliance on \texttt{gemini-3-flash-preview}
introduced a different class of limitation. While the model demonstrated strong
zero-shot pragmatic reasoning---evidenced by its capacity to identify sarcasm,
passive aggression, and culturally specific implicatures---its outputs were
probabilistic and not subject to human validation. The AI Partner's responses
were at times too lengthy and insufficiently naturalistic for fluid roleplay, and
the AI Mentor's scoring, while structurally enforced through JSON Schema
Controlled Generation, could not be guaranteed to be consistent across identical
inputs due to the stochastic nature of the model's generation process. Expert
evaluation of the chatbot's linguistic quality was not available to verify these
outputs independently.

Furthermore, the system's reliance on a single commercial model means that
findings cannot be generalized to alternative LLMs. Models with stronger
multilingual or code-switching capabilities, or those fine-tuned on Filipino
conversational data, may yield substantially different performance outcomes in
both the roleplay and pragmatic analysis functions.

\subsection{Testing Period and Sample}\label{subsec:test-limitations}

The evaluation design was constrained in ways that limit the strength of
conclusions that can be drawn about Socius's effectiveness in improving pragmatic
competence.

The testing period of one week was insufficient to capture meaningful or
statistically significant growth in pragmatic ability. Pragmatic competence---
encompassing the interpretation of social intent, cultural nuance, and
communicative strategy---is a complex, higher-order communicative skill that
develops gradually through repeated, varied, and reflective social practice.
Research in language education and social skill development consistently indicates
that such competencies require sustained exposure over weeks or months before
measurable behavioral change can be observed. The absence of quantitative
improvement in user scores over a five-day window therefore does not reflect a
failure of the system per se, but rather an evaluation window that is structurally
too short to detect the kind of longitudinal trajectory the system is designed to
support.

The sample was additionally limited in size and demographic scope. Of the 25
recruited participants, only 17 completed the full testing process, resulting in
an effective sample of 17. This level of attrition reduces statistical power and
narrows the basis for generalization. Furthermore, the sample was recruited through
convenience and purposive sampling, yielding a demographically homogeneous group
with limited coverage across age, regional background, and socioeconomic context.
While the study's theoretical grounding in Arnett's theory of Emerging Adulthood
defines the target population as broadly aged 18 to 24, the actual participant
demographics did not comprehensively represent the full range of this cohort,
particularly those from outside Metro Manila or from non-university contexts
where pragmatic challenges may manifest differently.

\subsection{Research Design}\label{subsec:design-limitations}

The evaluation of Socius was primarily focused on three methodological pillars:
functional testing of core system components, expert validation of the theoretical
framework and scenario quality, and end-user perception testing via Likert-based
evaluation forms. While these approaches provided meaningful insight into the
system's technical integrity and user experience, they also left several
substantive dimensions of the system's design and effectiveness unexamined.

Functional testing focused exclusively on component-level verification---
confirming that the conversation pipeline executed correctly, that the JSON
Schema scoring was produced and stored reliably, and that the Holistic Mentor
generated appropriate recommendations. What functional testing did not assess
was the alignment and compliance of the AI Partner's conversation flow with the
theoretical framework. Specifically, no systematic evaluation was conducted to
determine whether the AI Partner's in-character responses adhered to the
culturally grounded persona required by each scenario, whether Gricean maxim
framings were faithfully maintained across turns, or whether the pragmatic
complexity of the AI's dialogue was calibrated appropriately to the scenario's
intended developmental task.

Expert validation was conducted by a single evaluator with primary expertise in
the psychological dimensions of emerging adulthood and Filipino cultural values.
While this provided meaningful confirmation of the framework's structural soundness
and the scenarios' cultural relevance, it did not encompass the full domain
required for a comprehensive evaluation. In particular, no expert with specialization
in linguistics or pragmatics was engaged to evaluate the quality and validity of
the AI Mentor's pragmatic analysis outputs, the accuracy of the Gricean maxim
classifications, or the appropriateness of the feedback language generated for each
pragmatic category. The technical correctness of the chatbot's response quality
was likewise evaluated only by the research team, whose assessment was not
subject to independent review.

User testing, while valuable in capturing perceived usefulness and engagement,
suffered from several design limitations. AI Mentor usage was not made mandatory
during the testing period, which meant that a significant proportion of participants
did not engage with the feedback and analysis features, generating insufficient
data to evaluate the AI Mentor's effectiveness independently. No standardized
pre-test and post-test instrument for pragmatic ability was administered, meaning
there was no objective skills baseline against which to compare post-interaction
scores. Additionally, the presence of the research team during testing sessions,
combined with participants' awareness that their conversation logs were subject to
review, introduced observer bias and social desirability effects that may have
influenced the authenticity and risk-taking of users' conversational choices. The
text-only modality of the system further limited the scope of pragmatic dimensions
that could be practiced and observed, excluding auditory prosody, paralinguistic
cues, and non-verbal communication, all of which are integral to pragmatic
competence in Filipino interpersonal contexts.

%%%%%%%%%%%%%%%%%%%%%%%%%%%%%%%%%%%%%%%%%%%%%%%%%%
\section{Recommendations}\label{sec:recommendations}
%%%%%%%%%%%%%%%%%%%%%%%%%%%%%%%%%%%%%%%%%%%%%%%%%%

\subsection{Scientific and Technical Recommendations}\label{subsec:sci-tech-rec}

\paragraph{Resolving the Infrastructure Bottleneck.}
The most urgent technical priority is the remediation of infrastructure reliability.
User testing revealed that session disruptions attributable to Railway's free-tier
hosting eroded immersion and likely suppressed engagement with the AI Mentor's
feedback features. Future deployments should migrate to a paid cloud hosting tier
with guaranteed uptime and consistent compute availability---particularly for the
high-thinking-budget AI Mentor calls that are computationally expensive. Until
this is resolved, reliability constraints will continue to confound evaluations of
the system's pedagogical effectiveness, making it impossible to isolate whether
poor user engagement is caused by the content design or by the infrastructure.

\paragraph{Quantitative Validation of AI Mentor Scoring Consistency.}
Because the AI Mentor's pragmatic scoring is generated by a probabilistic LLM,
its inter-rater reliability and consistency have not been empirically established.
Future work should conduct a systematic inter-rater reliability study in which a
panel of human pragmatics experts independently scores a sample of conversation
transcripts using the same rubric criteria, and these human scores are compared
against the AI Mentor's outputs. Intraclass correlation coefficients or Cohen's
Kappa statistics should be computed per sub-dimension (Intention Coherence, Social
\& Affective Nuance, and Communicative Effectiveness) to determine where scoring
alignment is strong and where the model systematically diverges from expert
judgment. This is necessary before the scoring system can be treated as a valid
measure in future empirical studies.

\paragraph{Evaluating AI Partner Compliance with Theoretical Constructs.}
Functional testing confirmed that the conversation pipeline executes correctly but
did not assess whether the AI Partner's dialogue content was meaningfully aligned
with the culturally and theoretically grounded persona requirements of each
scenario. Future research should design a structured conversation-log analysis
protocol in which an expert evaluator---ideally with combined expertise in Filipino
linguistics and pragmatics---reviews a stratified sample of completed conversation
logs and assesses the following: (1) whether the AI Partner maintained its
designated cultural persona across turns, (2) whether Gricean maxim violations
were constructed intentionally and consistently within the scenario's design,
(3) whether the pragmatic complexity of the AI's responses was calibrated
appropriately to the scenario's developmental task, and (4) whether the AI
Mentor's feedback language was accurate, actionable, and grounded in the specific
conversational evidence from the transcript.

\paragraph{Exploring Alternative or Fine-Tuned Models.}
The system's current reliance on \texttt{gemini-3-flash-preview} means that its
performance characteristics are tied to a single commercial model. Future technical
work should explore whether models with stronger multilingual or code-switching
capabilities---particularly those trained on Tagalog-English conversational
corpora---produce more naturalistic AI Partner responses and more accurate
pragmatic analyses. The failed fine-tuning experiment with Mistral-7B documented
in this study established an important negative result: narrow dataset fine-tuning
on structured examples is insufficient for the complex, long-form reasoning
required by the AI Mentor. This finding argues for continued reliance on
generalist foundation models with sophisticated prompt engineering, rather than
fine-tuning smaller models on limited domain datasets.

\paragraph{Expanding Scenario Coverage and Pragmatic Breadth.}
User feedback identified that while existing scenarios were praised for cultural
relevance and realism, the AI Partner's responses sometimes failed to sustain
naturalistic conversational pressure across extended exchanges. This suggests that
the current scenario design, while theoretically sound, may not yet encode
sufficient variation in conversational dynamics---such as topic shifts, emotional
escalation, and repair sequences---to train users across the full pragmatic range
their scenarios imply. Future scenario development should introduce
sub-scenario branches that systematically vary the pragmatic difficulty, emotional
register, and repair demands of the dialogue, allowing users to practice adaptive
responses under more nuanced conditions.

\subsection{Social Science and Humanities Recommendations}\label{subsec:social-rec}

\paragraph{Extending the Evaluation Period to Capture Longitudinal Change.}
The absence of statistically significant quantitative change in users' pragmatic
scores over one week is best understood not as evidence that the system fails to
support learning, but as evidence that one week is structurally insufficient to
detect the kind of gradual, higher-order communicative development that Socius is
designed to foster. Pragmatic competence, like other complex social skills, develops
through accumulated practice, reflection, and transfer---processes that unfold over
weeks or months, not days. Future studies should implement a minimum four-to-six
week interaction period, with weekly or biweekly scenario sessions, to allow
sufficient exposure for measurable progress. A pre- and post-assessment using a
validated pragmatic competence instrument---distinct from the system's own
AI-generated scores---should be administered to provide an objective baseline
and endpoint measure against which system-tracked scores can be triangulated.

\paragraph{Expanding the Sample to Better Represent the Emerging Adulthood
Cohort.}
Arnett's theory of Emerging Adulthood defines this developmental stage as
broadly spanning ages 18 to 29, characterized by identity exploration, instability,
self-focus, and a feeling of being in-between. The current study's sample was
demographically narrow---concentrated among university-affiliated participants in
Metro Manila---which limits its ecological validity as a test of a system intended
for the wider Filipino emerging adult population. Future studies should recruit
participants across the full 18-to-29 age range and include individuals from diverse
regional, socioeconomic, and educational backgrounds. This expansion is warranted
not merely for statistical power, but because pragmatic norms, cultural values,
and social developmental pressures are known to vary significantly across these
demographic dimensions in the Philippine context. A more representative sample
would allow the study to assess whether the system's cultural grounding generalizes
beyond the narrowly defined participant group tested in this iteration.

\paragraph{Incorporating Mandatory AI Mentor Engagement and Post-Session
Reflection.}
A critical design gap in the current evaluation was that AI Mentor usage was
optional, resulting in insufficient data to assess its effectiveness as a feedback
instrument. This limitation is not merely a research design issue but a pedagogical
one: structured, timely feedback is one of the most consistently supported
predictors of skill development in social and communicative training contexts.
Future iterations should make AI Mentor engagement a required component of each
session, and should additionally introduce a brief structured reflection prompt
at the end of each conversation in which users articulate one pragmatic insight
they gained and one communicative strategy they intend to apply differently in
future interactions. This active retrieval and reflection step has strong empirical
support in the learning sciences and would serve both pedagogical and evaluative
purposes by generating richer qualitative data on perceived learning.

\paragraph{Addressing Psychological and Social Risks of AI-Mediated Social
Practice.}
The finding that users felt emotionally invested in the scenarios---to the point
of describing the experience as ``eerie'' and realistic---raises important
questions about the psychological and social risks of sustained AI-mediated
social skills training. On one hand, this emotional engagement is a marker of
the system's effectiveness in simulating authentic conversational pressure. On
the other hand, users also acknowledged that their awareness of talking to an AI
may have led them to take fewer communicative risks than they would in real social
encounters, reducing the authenticity of their practice. Future research should
engage a psychologist or social scientist to conduct a post-user-conversation-log
analysis, evaluating the following: (1) whether users demonstrated safe, avoidant,
or risk-averse communicative strategies that would not serve them in real
interactions; (2) whether repeated exposure to AI-generated social conflict
scenarios produces adaptive or maladaptive pragmatic patterns; (3) whether the
safety of the AI context transferred into greater real-world communicative
confidence, or whether it reinforced a preference for low-stakes digital
interaction over authentic social engagement; and (4) whether the AI Partner's
capacity to simulate toxic social dynamics---such as deflection, passive
aggression, and boundary violations---poses any psychological risks for users
with histories of interpersonal conflict or social anxiety.

\paragraph{Toward Real-World Transfer: Bridging Simulation and Social Practice.}
A recurring theme in user feedback was the gap between performance within Socius
and transfer of skills to real-world social interactions. While the safety of the
AI environment supported user engagement, it simultaneously reduced the stakes
that drive genuine pragmatic adaptation. Future versions of the system should
incorporate explicit transfer-bridging mechanisms---such as post-session reflection
prompts that ask users to identify an upcoming real-world scenario analogous to
the one they just practiced, or structured debriefing conversations with the AI
Mentor that focus not only on what the user did during the session but on how
they would apply those insights in an actual social encounter. The effectiveness
of such bridging strategies should itself be evaluated in subsequent studies,
as the gap between AI-assisted rehearsal and authentic social competence
represents one of the most significant unresolved challenges in conversational AI
coaching research.